\chapter{Related Work}
\label{ch:related}

The literature relevant to this project falls into four groups: classical approaches to estimating network performance, learned network models (with GNN-based models as the main line), learning-based routing in SDN, and model compression with an emphasis on knowledge distillation for GNNs. The chapter closes with a comparison across these groups and a statement of the research gap.

\section{Analytical and Simulation-Based Performance Models}
\label{sec:rw-classic}

Queueing theory gives closed-form or numerically cheap estimates of delay and loss for individual queues and, under product-form assumptions, for networks of queues \cite{kleinrock1975queueing}. Its strength is speed and interpretability. Its weakness is that the assumptions needed for tractability (Markovian arrivals, independence between queues, simple scheduling) are routinely broken in practice. Bursty or autocorrelated traffic, priority and weighted-fair scheduling, and correlation between hops along a path all push measured delay away from the analytical estimate. Ferriol-Galm\'es et~al.\ \cite{ferriolgalmes2023routenetfermi} include a queueing-theory baseline in their evaluation and find it considerably less accurate than their learned model once traffic departs from the Poisson assumption.

Discrete-event packet simulators, such as OMNeT++ \cite{varga2001omnet} and ns-3 \cite{riley2010ns3}, model every packet and so capture these effects faithfully. The price is run time, which grows with the number of packets and therefore with network size and load. Simulators are the usual source of ground truth for training learned models, including the dataset used in this project, but they are too slow to sit inside a controller's decision loop. Kim et~al.\ \cite{kim2022drlsdn} make a similar observation from the routing side: they train a DRL agent against an offline M/M/1/K queueing model precisely because interacting with a live network or a full simulator during training would be too slow and too disruptive.

The two families therefore occupy opposite ends of a speed--accuracy spectrum, and the motivation for learned models is to occupy the middle of it.

\section{Learned Network Models}
\label{sec:rw-learned}

\subsection{Non-graph models}

Early learned network models treated the network as a fixed-size input vector. The knowledge-defined networking proposal \cite{mestres2017kdn} trained a multi-layer perceptron to map a traffic matrix to delay for one specific topology. Such models do not transfer to a different topology, because the input layout is tied to the topology on which they were trained. Later systems kept accuracy but changed the abstraction. MimicNet \cite{zhang2021mimicnet} learns per-cluster models of a data-centre network and composes them to estimate the behaviour of a much larger fabric, and DeepQueueNet \cite{yang2022deepqueuenet} learns packet-level device models that are assembled into a network-wide simulation. Both scale better than packet simulation, but both are designed around data-centre structure or around replaying packet traces, which is a different setting from per-flow delay prediction for arbitrary wide-area topologies.

\subsection{Graph neural network models}

GNNs \cite{scarselli2009gnn,gilmer2017mpnn,battaglia2018relational} compute a hidden state for every entity in a graph by repeatedly exchanging messages with neighbouring entities. Because the same learned functions are applied to every node and edge, a trained GNN can be run on graphs of a different size and shape from those it saw in training, which is the property that network modelling needs.

RouteNet \cite{rusek2020routenet} introduced a GNN tailored to networks. It represents links and paths as two sets of entities and alternates between updating each path from the links it traverses (using a recurrent unit over the ordered link sequence) and updating each link from the paths that cross it. Trained on simulated data, RouteNet predicted per-path delay and jitter on topologies outside its training set and was used inside an optimisation loop to choose routings that meet delay targets. Its main limitation was that it assumed a single first-in, first-out queue per link and Poisson-like traffic.

RouteNet-Erlang \cite{ferriolgalmes2022routeneterlang} and then RouteNet-Fermi \cite{ferriolgalmes2023routenetfermi} removed these limitations. RouteNet-Fermi adds queues as a third entity type, so that a link can host several queues with different priorities and weights, and supports four scheduling policies (FIFO, strict priority, WFQ and DRR) and several traffic models, including on--off and autoregressive sources. Rather than outputting delay directly, it predicts per-hop queue occupancy and converts it to delay using the link capacity, which makes the output scale correctly when link speeds change. The authors report that this design generalises to networks larger than those used for training. RouteNet-Fermi is the teacher architecture in this project and is described in detail in \cref{sec:teacher}.

A shared feature of this line of work is that model size and computational cost are not the main concern. RouteNet-Fermi uses three gated recurrent unit (GRU) cells \cite{cho2014gru}, 32-dimensional states and eight message-passing iterations. Jiang et~al.\ \cite{jiang2024gnnrouting} identify scalability, inference latency and memory use as open problems for GNNs in routing optimisation, and the RouteNet papers themselves evaluate accuracy and generalisation rather than deployment cost on constrained hardware.

\section{Learning-Based Routing in SDN}
\label{sec:rw-routing}

A related body of work uses learning to make routing decisions directly. Kim et~al.\ \cite{kim2022drlsdn} train a DRL agent to set link weights so as to reduce end-to-end delay and loss in SDN. Song et~al.\ \cite{song2024qosdrl} use deep deterministic policy gradient with a reward that combines delay, bandwidth, jitter and loss. Almasan et~al.\ \cite{almasan2022drlgnn} show that a DRL agent built on a GNN generalises across topologies far better than one built on a fully connected network, because the GNN does not tie its input to a fixed topology. He et~al.\ \cite{he2024causality} add causality detection to a GNN-based DRL agent to make it more stable under noisy conditions, and Ampratwum and Nayak \cite{ampratwum2025wdm} combine DRL and GNNs for restoration in optical networks.

These works matter to this project in two ways. First, several of them depend, during training or at decision time, on a model that estimates the consequence of an action; a fast and accurate delay model is therefore a building block for this kind of routing. Second, their evaluations generally measure routing quality rather than inference cost, and those that use GNNs inherit the computational load of the underlying GNN. Neither line of work, as reviewed here, studies how small the network model can be made before its predictions become unreliable.

\section{Model Compression and Knowledge Distillation}
\label{sec:rw-kd}

\subsection{General compression techniques}

Compression methods reduce the cost of running a trained network. Pruning removes weights or structures that contribute little \cite{han2016deepcompression}, and quantisation stores and computes with low-precision integers in place of 32-bit floats \cite{jacob2018quantization}. Both start from a trained architecture and make it cheaper while keeping its structure. Knowledge distillation \cite{hinton2015distilling} instead trains a separate, smaller model. In its original form the student is trained on a weighted sum of the usual loss against ground-truth labels and a loss that matches the teacher's softened output distribution; later variants also match intermediate representations \cite{romero2015fitnets}. A survey by Gou et~al.\ \cite{gou2021kdsurvey} classifies the many variants by what knowledge is transferred (responses, features or relations) and how (offline, online or self-distillation). Distillation is attractive when the desired student differs structurally from the teacher, which is the case in this project.

\subsection{Distillation for graph neural networks}

Applying distillation to GNNs raises the question of what graph-specific knowledge should be transferred. Yang et~al.\ \cite{yang2020lsp} propose a local-structure-preserving module that makes the student reproduce the teacher's view of how each node relates to its neighbours, not only its final outputs. Zhang et~al.\ \cite{zhang2022glnn} distil a GNN into a plain MLP that uses no graph structure at inference time; the resulting student is much faster but can only rely on the features of each node, which suits node classification on a fixed graph. Chen et~al.\ \cite{chen2021gnnsd} remove the separate teacher altogether and propose self-distillation, in which a GNN regularises its own deeper layers using a measure of how much each node's embedding differs from its neighbourhood. Yang et~al.\ \cite{yang2023ltd} learn distillation temperatures for individual nodes jointly with the student, rather than using one global temperature. Huo et~al.\ \cite{huo2023t2gnn} use two teachers, one specialised in node features and one in graph structure, to train a student that copes with graphs whose features and edges are partly missing. LightGNN \cite{chen2025lightgnn} combines pruning with hierarchical distillation to compress GNN-based recommender models.

Two features of this work are worth noting for the present project. First, almost all of it concerns classification or ranking: the teacher's output is a probability distribution or a score, and temperature-softened outputs carry ``dark knowledge'' about relative class similarity. Per-flow delay is a continuous regression target with no such distribution, so the classical soft-label loss does not apply directly and a regression loss on the teacher's predictions has to be used instead. Second, the student architectures studied are usually smaller copies of the teacher, or MLPs that drop the graph entirely. A network model cannot drop the graph, because the delay of a flow depends on what other flows share its links and queues, so a student for this task has to remain a GNN while still being cheaper.

\section{Synthesis and Research Gap}
\label{sec:gap}

\Cref{tab:rw-compare} compares representative work along the dimensions that matter for this project.

\begin{table}[ht]
\centering
\small
\caption{Comparison of representative related work against the requirements of this project.}
\label{tab:rw-compare}
\begin{tabularx}{\textwidth}{@{}l Y C{19mm} C{19mm} C{17mm}@{}}
\toprule
\textbf{Approach} & \textbf{Representative work} & \textbf{Per-flow QoS} & \textbf{Unseen topologies} & \textbf{Low cost} \\
\midrule
Queueing theory      & \cite{kleinrock1975queueing}                       & Approx. & Yes & Yes \\
Packet simulation    & \cite{varga2001omnet,riley2010ns3}                 & Yes     & Yes & No  \\
Non-graph learned    & \cite{mestres2017kdn,zhang2021mimicnet,yang2022deepqueuenet} & Partly & Limited & Partly \\
GNN network models   & \cite{rusek2020routenet,ferriolgalmes2022routeneterlang,ferriolgalmes2023routenetfermi} & Yes & Yes & Not studied \\
DRL / DRL+GNN routing & \cite{kim2022drlsdn,song2024qosdrl,almasan2022drlgnn,he2024causality} & Decision, not estimate & Partly & Not studied \\
GNN distillation     & \cite{yang2020lsp,zhang2022glnn,chen2021gnnsd,yang2023ltd,huo2023t2gnn} & Not addressed & n/a & Yes \\
\bottomrule
\end{tabularx}
\end{table}

The comparison shows a clear gap. GNN network models such as RouteNet-Fermi provide accurate per-flow estimates that transfer across topologies, but their computational cost on constrained hardware has not been the subject of study. GNN distillation methods reduce cost, but they have been developed and evaluated for classification and recommendation, with students that are either smaller copies of the teacher or graph-free MLPs. No reviewed work has examined whether a structurally simpler GNN student, distilled from a network-modelling teacher, can keep useful delay accuracy, nor how such a student behaves on topologies larger than those seen in training. A related point, which emerges from the experiments in \cref{ch:results}, is that the choice of training topologies for the teacher itself has a large effect on this out-of-distribution accuracy.

This project addresses that gap by (i) measuring the effect of topology diversity on RouteNet-Fermi's accuracy on larger, unseen topologies, (ii) proposing a student architecture that merges queues into links to shorten message passing, and (iii) evaluating the student with and without distillation on the same held-out topologies.
